\section{Quantitative Analysis}
\label{sec:experiments_latent_space_analysis}

This section analyzes the \gls{ls} of the surrogate policies using \gls{cu}, \gls{as}, and \gls{acu}.
For each environment and each layer, a random subset of 20,000 observations is embedded into the \gls{ls} of each surrogate policy, and the resulting point cloud $\textcolor{myblue}{\mathcal{E}}$ is partitioned using $k$-means clustering \cite{macqueen1967multivariate} for each number of clusters $K \in [2,20]$.
The main results are presented in this section, while the full breakdown across all analyzed layers is reported in Appendix~\ref{chap}.

Figure~\ref{fig:latent_space_curves_by_layer_depth} reports, for each environment, \gls{cu}, \gls{as}, and \gls{acu} averaged over $K \in [2,20]$ with respect to the depth of the analyzed layer.
In every environment, \gls{cu} tends to be highest in the earliest layers and to decrease with depth, while \gls{as} follows the opposite trend, increasing with depth as the representation becomes less tied to the raw observation, the action distribution, and the network's initialization.
\gls{acu}, which combines both effects, therefore increases with depth for \MarioBrosName, \MsPacManName, and \SurroundName, for which the last analyzed layer is the most informative, whereas it peaks at an intermediate depth and decreases again for \BreakoutName, \PongName, and \SpaceInvadersName.

% Spacing specific to this figure (tuned for its own subfigure heights --
% do not share these with other multi-subfigure figures, see
% figures/clustering_agent_representation/latent_space_curves_by_metric.tex
% and figures/experiments/policy_training.tex for those figures' own copies).
\newcommand{\LatentCurvesByDepthSubfigCaptionSkip}{0.1cm} % space between a subfigure's plot and its own caption
\newcommand{\LatentCurvesByDepthSubfigSpacing}{-0.2cm}   % vertical space between consecutive subfigures
\newcommand{\LatentCurvesByDepthSubfigLegendSpacing}{0.2cm} % vertical space between the last subfigure and the shared legend

\begin{figure}[t]
    \centering
    \begin{subfigure}[b]{\textwidth}
        \centering
        \captionsetup{skip=\LatentCurvesByDepthSubfigCaptionSkip}
        \resizebox{\linewidth}{!}{
        \begin{tikzpicture}
        \begin{axis}[
            xlabel={Layer depth},
            ylabel={\gls{cu}},
            grid=major,
            xmin=1, xmax=8,
            xtick distance=1,
            ymin=0, ymax=1,
            ytick distance=0.25,
            ymajorgrids=true,
            width=\linewidth,
            height=0.33\linewidth,
        ]
            \addplot[color=breakoutColor, mark=\breakoutMark, mark size=\breakoutMarkSize] coordinates {
                (1,0.85) (2,0.73) (3,0.57) (4,0.50) (5,0.50) (6,0.59) (7,0.50) (8,0.51)
            };
            \addplot[color=marioBrosColor, mark=\marioBrosMark, mark size=\marioBrosMarkSize] coordinates {
                (1,0.62) (2,0.68) (3,0.62) (4,0.55) (5,0.55) (6,0.48) (7,0.40) (8,0.37)
            };
            \addplot[color=msPacmanColor, mark=\msPacmanMark, mark size=\msPacmanMarkSize] coordinates {
                (1,0.78) (2,0.72) (3,0.68) (4,0.70) (5,0.70) (6,0.58) (7,0.51) (8,0.58)
            };
            \addplot[color=pongColor, mark=\pongMark, mark size=\pongMarkSize] coordinates {
                (1,0.62) (2,0.65) (3,0.48) (4,0.50) (5,0.50) (6,0.31)
            };
            \addplot[color=spaceInvadersColor, mark=\spaceInvadersMark, mark size=\spaceInvadersMarkSize] coordinates {
                (1,0.75) (2,0.81) (3,0.74) (4,0.63) (5,0.63) (6,0.49)
            };
            \addplot[color=surroundColor, mark=\surroundMark, mark size=\surroundMarkSize] coordinates {
                (1,0.70) (2,0.58) (3,0.66) (4,0.65) (5,0.65) (6,0.59) (7,0.49) (8,0.54)
            };
        \end{axis}
        \end{tikzpicture}
        }
        \caption{\gls{cu} with respect to the layer depth.}
    \end{subfigure}
    \vspace{\LatentCurvesByDepthSubfigSpacing}

    \begin{subfigure}[b]{\textwidth}
        \centering
        \captionsetup{skip=\LatentCurvesByDepthSubfigCaptionSkip}
        \resizebox{\linewidth}{!}{
        \begin{tikzpicture}
        \begin{axis}[
            xlabel={Layer depth},
            ylabel={\gls{as}},
            grid=major,
            xmin=1, xmax=8,
            xtick distance=1,
            ymin=0, ymax=1,
            ytick distance=0.25,
            ymajorgrids=true,
            width=\linewidth,
            height=0.33\linewidth,
        ]
            \addplot[color=breakoutColor, mark=\breakoutMark, mark size=\breakoutMarkSize] coordinates {
                (1,0.15) (2,0.28) (3,0.44) (4,0.76) (5,0.76) (6,0.82) (7,0.79) (8,0.76)
            };
            \addplot[color=marioBrosColor, mark=\marioBrosMark, mark size=\marioBrosMarkSize] coordinates {
                (1,0.44) (2,0.54) (3,0.62) (4,0.66) (5,0.66) (6,0.71) (7,0.87) (8,0.94)
            };
            \addplot[color=msPacmanColor, mark=\msPacmanMark, mark size=\msPacmanMarkSize] coordinates {
                (1,0.19) (2,0.27) (3,0.35) (4,0.41) (5,0.41) (6,0.55) (7,0.64) (8,0.74)
            };
            \addplot[color=pongColor, mark=\pongMark, mark size=\pongMarkSize] coordinates {
                (1,0.33) (2,0.59) (3,0.91) (4,0.95) (5,0.95) (6,0.95)
            };
            \addplot[color=spaceInvadersColor, mark=\spaceInvadersMark, mark size=\spaceInvadersMarkSize] coordinates {
                (1,0.28) (2,0.45) (3,0.54) (4,0.72) (5,0.72) (6,0.74)
            };
            \addplot[color=surroundColor, mark=\surroundMark, mark size=\surroundMarkSize] coordinates {
                (1,0.27) (2,0.45) (3,0.56) (4,0.49) (5,0.49) (6,0.57) (7,0.72) (8,0.77)
            };
        \end{axis}
        \end{tikzpicture}
        }
        \caption{\gls{as} with respect to the layer depth.}
    \end{subfigure}
    \vspace{\LatentCurvesByDepthSubfigSpacing}

    \begin{subfigure}[b]{\textwidth}
        \centering
        \captionsetup{skip=\LatentCurvesByDepthSubfigCaptionSkip}
        \resizebox{\linewidth}{!}{
        \begin{tikzpicture}
        \begin{axis}[
            xlabel={Layer depth},
            ylabel={\gls{acu}},
            grid=major,
            xmin=1, xmax=8,
            xtick distance=1,
            ymin=-0.2, ymax=1,
            ytick distance=0.2,
            ymajorgrids=true,
            width=\linewidth,
            height=0.33\linewidth,
        ]
            \addplot[color=breakoutColor, mark=\breakoutMark, mark size=\breakoutMarkSize] coordinates {
                (1,0.00) (2,0.01) (3,0.01) (4,0.27) (5,0.27) (6,0.41) (7,0.29) (8,0.27)
            };
            \addplot[color=marioBrosColor, mark=\marioBrosMark, mark size=\marioBrosMarkSize] coordinates {
                (1,0.06) (2,0.22) (3,0.24) (4,0.21) (5,0.21) (6,0.19) (7,0.27) (8,0.31)
            };
            \addplot[color=msPacmanColor, mark=\msPacmanMark, mark size=\msPacmanMarkSize] coordinates {
                (1,-0.04) (2,-0.01) (3,0.04) (4,0.10) (5,0.10) (6,0.14) (7,0.15) (8,0.32)
            };
            \addplot[color=pongColor, mark=\pongMark, mark size=\pongMarkSize] coordinates {
                (1,-0.05) (2,0.24) (3,0.39) (4,0.45) (5,0.45) (6,0.27)
            };
            \addplot[color=spaceInvadersColor, mark=\spaceInvadersMark, mark size=\spaceInvadersMarkSize] coordinates {
                (1,0.03) (2,0.25) (3,0.28) (4,0.35) (5,0.35) (6,0.24)
            };
            \addplot[color=surroundColor, mark=\surroundMark, mark size=\surroundMarkSize] coordinates {
                (1,-0.02) (2,0.03) (3,0.22) (4,0.14) (5,0.14) (6,0.16) (7,0.21) (8,0.31)
            };
        \end{axis}
        \end{tikzpicture}
        }
        \caption{\gls{acu} with respect to the layer depth.}
    \end{subfigure}

    \vspace{\LatentCurvesByDepthSubfigLegendSpacing}
    {\hypersetup{pdfborder={0 0 0}}\ref{policyTrainingLegend}}

    \caption{\gls{cu}, \gls{as}, and \gls{acu} averaged over the number of clusters $K \in [2,20]$, with respect to the depth of the analyzed layer, for each of the six environments.}
    \label{fig:latent_space_curves_by_layer_depth}
\end{figure}

% (see Appendix~\ref{chap:appendix} for the full per-layer, per-$K$ values underlying this figure)


The linear fit of the environment-average curve confirms these trends: \gls{cu} decreases roughly linearly with depth, from $0.72$ at the first layer to $0.50$ at the last (slope $\approx -0.04$ per layer), while \gls{as} increases at a comparable rate in the opposite direction, from $0.28$ to $0.80$ (slope $\approx +0.07$ per layer).
\gls{acu} increases more moderately on average, from $0.00$ to $0.30$ (slope $\approx +0.03$ per layer), indicating that, despite the environment-specific peaks and decreases described above, deeper representations tend on average to offer a better trade-off between \gls{cu} and \gls{as}.

For each environment, the results are reported for the layer that achieves the highest average \gls{acu} across the tested values of $K$.
This layer is therefore the one for which the agreement between independently trained surrogate policies is least explained by the observation space, the action space, or the network initialization.
Table~\ref{tab:latent_space_best_layer} reports this layer for each environment, together with the corresponding average \gls{cu}, \gls{as}, and \gls{acu} over $K$; the evolution of these three metrics with respect to $K$ for that layer is then shown in Figure~\ref{fig:latent_space_curves_by_metric}.

\begin{table}[t]
\centering
\resizebox{0.8\linewidth}{!}{
    \begin{tblr}{
        colspec={|c|c|c|c|c|},
        hlines,
        vlines,
        row{1} = {font=\bfseries, halign=c, valign=m}
    }
        Environment     & Layer & \gls{cu}        & \gls{as}        & \gls{acu}       \\
        Breakout        & 6b    & 0.59 $\pm$0.06  & 0.82 $\pm$0.05  & 0.41 $\pm$0.06  \\
        Mario Bros      & 8b    & 0.37 $\pm$0.08  & 0.94 $\pm$0.02  & 0.31 $\pm$0.08  \\
        Ms.\ Pac-Man    & 8b    & 0.58 $\pm$0.08  & 0.74 $\pm$0.05  & 0.32 $\pm$0.07  \\
        Pong            & 4     & 0.50 $\pm$0.04  & 0.95 $\pm$0.03  & 0.45 $\pm$0.07  \\
        Space Invaders  & 4     & 0.63 $\pm$0.11  & 0.73 $\pm$0.06  & 0.35 $\pm$0.09  \\
        Surround        & 8b    & 0.54 $\pm$0.10  & 0.77 $\pm$0.07  & 0.31 $\pm$0.06  \\
    \end{tblr}
}
\caption{Layer that maximizes the average \gls{acu} across $K \in [2,20]$ for each environment, together with the corresponding average \gls{cu}, \gls{as}, and \gls{acu} over $K$.}
\label{tab:latent_space_best_layer}
\end{table}


% Spacing specific to this figure (tuned for its own subfigure heights --
% do not share these with other multi-subfigure figures, see
% figures/experiments/policy_training.tex for that figure's own copies).
\newcommand{\LatentCurvesSubfigCaptionSkip}{0.1cm} % space between a subfigure's plot and its own caption
\newcommand{\LatentCurvesSubfigSpacing}{-0.2cm}   % vertical space between consecutive subfigures
\newcommand{\LatentCurvesSubfigLegendSpacing}{0.2cm} % vertical space between the last subfigure and the shared legend
\newcommand{\LatentCurvesSubfigHeight}{0.33\linewidth} % height of each subfigure's plot (controls this figure's overall scale)

\begin{figure}[H]
    \centering
    \begin{subfigure}[b]{\textwidth}
        \centering
        \captionsetup{skip=\LatentCurvesSubfigCaptionSkip}
        \resizebox{\linewidth}{!}{
        \begin{tikzpicture}
        \begin{axis}[
            xlabel={Number of clusters},
            ylabel={\gls{cu}},
            grid=major,
            xmin=1, xmax=20,
            ymin=0, ymax=1,
            ytick distance=0.25,
            ymajorgrids=true,
            width=\linewidth,
            height=\LatentCurvesSubfigHeight,
        ]
            \addplot[color=breakoutColor, mark=\breakoutMark, mark size=\breakoutMarkSize] coordinates {
                (2,0.66) (3,0.57) (4,0.65) (5,0.79) (6,0.59) (7,0.62) (8,0.56) (9,0.56) (10,0.58) (11,0.64) (12,0.59) (13,0.55) (14,0.55) (15,0.54) (16,0.59) (17,0.60) (18,0.53) (19,0.55) (20,0.55)
            };
            \addplot[color=marioBrosColor, mark=\marioBrosMark, mark size=\marioBrosMarkSize] coordinates {
                (2,0.99) (3,0.84) (4,0.58) (5,0.75) (6,0.54) (7,0.57) (8,0.57) (9,0.66) (10,0.62) (11,0.62) (12,0.58) (13,0.62) (14,0.59) (15,0.55) (16,0.50) (17,0.56) (18,0.59) (19,0.54) (20,0.52)
            };
            \addplot[color=msPacmanColor, mark=\msPacmanMark, mark size=\msPacmanMarkSize] coordinates {
                (2,0.66) (3,0.78) (4,0.58) (5,0.60) (6,0.67) (7,0.64) (8,0.64) (9,0.62) (10,0.59) (11,0.59) (12,0.53) (13,0.62) (14,0.61) (15,0.48) (16,0.48) (17,0.45) (18,0.49) (19,0.48) (20,0.46)
            };
            \addplot[color=pongColor, mark=\pongMark, mark size=\pongMarkSize] coordinates {
                (2,0.60) (3,0.52) (4,0.58) (5,0.49) (6,0.44) (7,0.51) (8,0.58) (9,0.48) (10,0.50) (11,0.52) (12,0.49) (13,0.46) (14,0.50) (15,0.45) (16,0.46) (17,0.47) (18,0.48) (19,0.47) (20,0.53)
            };
            \addplot[color=spaceInvadersColor, mark=\spaceInvadersMark, mark size=\spaceInvadersMarkSize] coordinates {
                (2,0.31) (3,0.73) (4,0.92) (5,0.69) (6,0.61) (7,0.57) (8,0.64) (9,0.61) (10,0.61) (11,0.66) (12,0.58) (13,0.52) (14,0.56) (15,0.63) (16,0.67) (17,0.63) (18,0.63) (19,0.68) (20,0.64)
            };
            \addplot[color=surroundColor, mark=\surroundMark, mark size=\surroundMarkSize] coordinates {
                (2,0.77) (3,0.75) (4,0.67) (5,0.54) (6,0.52) (7,0.55) (8,0.52) (9,0.58) (10,0.61) (11,0.59) (12,0.51) (13,0.52) (14,0.50) (15,0.46) (16,0.50) (17,0.45) (18,0.43) (19,0.43) (20,0.43)
            };
        \end{axis}
        \end{tikzpicture}
        }
        \caption{\gls{cu} with respect to the number of clusters.}
    \end{subfigure}
    \vspace{\LatentCurvesSubfigSpacing}

    \begin{subfigure}[b]{\textwidth}
        \centering
        \captionsetup{skip=\LatentCurvesSubfigCaptionSkip}
        \resizebox{\linewidth}{!}{
        \begin{tikzpicture}
        \begin{axis}[
            xlabel={Number of clusters},
            ylabel={\gls{as}},
            grid=major,
            xmin=1, xmax=20,
            ymin=0, ymax=1,
            ytick distance=0.25,
            ymajorgrids=true,
            width=\linewidth,
            height=\LatentCurvesSubfigHeight,
        ]
            \addplot[color=breakoutColor, mark=\breakoutMark, mark size=\breakoutMarkSize] coordinates {
                (2,0.62) (3,0.83) (4,0.74) (5,0.82) (6,0.85) (7,0.85) (8,0.85) (9,0.85) (10,0.83) (11,0.82) (12,0.85) (13,0.85) (14,0.84) (15,0.83) (16,0.83) (17,0.83) (18,0.84) (19,0.85) (20,0.83)
            };
            \addplot[color=marioBrosColor, mark=\marioBrosMark, mark size=\marioBrosMarkSize] coordinates {
                (2,0.00) (3,0.66) (4,0.63) (5,0.65) (6,0.64) (7,0.62) (8,0.63) (9,0.62) (10,0.69) (11,0.69) (12,0.64) (13,0.66) (14,0.62) (15,0.64) (16,0.67) (17,0.65) (18,0.66) (19,0.64) (20,0.69)
            };
            \addplot[color=msPacmanColor, mark=\msPacmanMark, mark size=\msPacmanMarkSize] coordinates {
                (2,0.56) (3,0.59) (4,0.47) (5,0.57) (6,0.56) (7,0.54) (8,0.55) (9,0.56) (10,0.60) (11,0.58) (12,0.62) (13,0.59) (14,0.59) (15,0.65) (16,0.65) (17,0.68) (18,0.66) (19,0.66) (20,0.67)
            };
            \addplot[color=pongColor, mark=\pongMark, mark size=\pongMarkSize] coordinates {
                (2,0.99) (3,0.99) (4,0.98) (5,0.99) (6,0.97) (7,0.96) (8,0.97) (9,0.96) (10,0.93) (11,0.94) (12,0.96) (13,0.93) (14,0.93) (15,0.93) (16,0.93) (17,0.91) (18,0.90) (19,0.90) (20,0.91)
            };
            \addplot[color=spaceInvadersColor, mark=\spaceInvadersMark, mark size=\spaceInvadersMarkSize] coordinates {
                (2,0.94) (3,0.80) (4,0.69) (5,0.71) (6,0.71) (7,0.75) (8,0.74) (9,0.74) (10,0.75) (11,0.69) (12,0.67) (13,0.68) (14,0.71) (15,0.68) (16,0.68) (17,0.70) (18,0.69) (19,0.74) (20,0.69)
            };
            \addplot[color=surroundColor, mark=\surroundMark, mark size=\surroundMarkSize] coordinates {
                (2,0.66) (3,0.53) (4,0.71) (5,0.70) (6,0.71) (7,0.70) (8,0.71) (9,0.69) (10,0.68) (11,0.71) (12,0.74) (13,0.74) (14,0.73) (15,0.74) (16,0.73) (17,0.75) (18,0.76) (19,0.75) (20,0.76)
            };
        \end{axis}
        \end{tikzpicture}
        }
        \caption{\gls{as} with respect to the number of clusters.}
    \end{subfigure}
    \vspace{\LatentCurvesSubfigSpacing}

    \begin{subfigure}[b]{\textwidth}
        \centering
        \captionsetup{skip=\LatentCurvesSubfigCaptionSkip}
        \resizebox{\linewidth}{!}{
        \begin{tikzpicture}
        \begin{axis}[
            xlabel={Number of clusters},
            ylabel={\gls{acu}},
            grid=major,
            xmin=1, xmax=20,
            ymin=-0.2, ymax=1,
            ytick distance=0.2,
            ymajorgrids=true,
            width=\linewidth,
            height=\LatentCurvesSubfigHeight,
        ]
            \addplot[color=breakoutColor, mark=\breakoutMark, mark size=\breakoutMarkSize] coordinates {
                (2,0.28) (3,0.40) (4,0.39) (5,0.61) (6,0.43) (7,0.47) (8,0.41) (9,0.41) (10,0.42) (11,0.46) (12,0.45) (13,0.40) (14,0.39) (15,0.37) (16,0.42) (17,0.43) (18,0.36) (19,0.40) (20,0.38)
            };
            \addplot[color=marioBrosColor, mark=\marioBrosMark, mark size=\marioBrosMarkSize] coordinates {
                (2,0.00) (3,0.50) (4,0.21) (5,0.40) (6,0.18) (7,0.20) (8,0.20) (9,0.27) (10,0.31) (11,0.31) (12,0.23) (13,0.28) (14,0.21) (15,0.19) (16,0.16) (17,0.22) (18,0.25) (19,0.18) (20,0.21)
            };
            \addplot[color=msPacmanColor, mark=\msPacmanMark, mark size=\msPacmanMarkSize] coordinates {
                (2,0.22) (3,0.37) (4,0.05) (5,0.16) (6,0.23) (7,0.18) (8,0.19) (9,0.19) (10,0.19) (11,0.17) (12,0.15) (13,0.21) (14,0.20) (15,0.13) (16,0.12) (17,0.13) (18,0.15) (19,0.14) (20,0.13)
            };
            \addplot[color=pongColor, mark=\pongMark, mark size=\pongMarkSize] coordinates {
                (2,0.59) (3,0.51) (4,0.56) (5,0.48) (6,0.41) (7,0.48) (8,0.56) (9,0.44) (10,0.44) (11,0.46) (12,0.45) (13,0.40) (14,0.43) (15,0.38) (16,0.39) (17,0.38) (18,0.38) (19,0.38) (20,0.44)
            };
            \addplot[color=spaceInvadersColor, mark=\spaceInvadersMark, mark size=\spaceInvadersMarkSize] coordinates {
                (2,0.26) (3,0.53) (4,0.61) (5,0.40) (6,0.32) (7,0.32) (8,0.38) (9,0.36) (10,0.35) (11,0.35) (12,0.26) (13,0.20) (14,0.27) (15,0.31) (16,0.36) (17,0.33) (18,0.32) (19,0.43) (20,0.34)
            };
            \addplot[color=surroundColor, mark=\surroundMark, mark size=\surroundMarkSize] coordinates {
                (2,0.42) (3,0.28) (4,0.38) (5,0.24) (6,0.23) (7,0.24) (8,0.23) (9,0.27) (10,0.29) (11,0.30) (12,0.24) (13,0.26) (14,0.23) (15,0.20) (16,0.23) (17,0.20) (18,0.18) (19,0.18) (20,0.19)
            };
        \end{axis}
        \end{tikzpicture}
        }
        \caption{\gls{acu} with respect to the number of clusters.}
    \end{subfigure}

    \vspace{\LatentCurvesSubfigLegendSpacing}
    {\hypersetup{pdfborder={0 0 0}}\ref{policyTrainingLegend}}

    \caption{\gls{cu}, \gls{as}, and \gls{acu} with respect to the number of clusters, for each of the six environments, for the layer that maximizes the average \gls{acu} in that environment.}
    \label{fig:latent_space_curves_by_metric}
\end{figure}

% (see Table~\ref{tab:latent_space_best_layer} for the specific layer and average values of each environment)

The best layer is the last analyzed one for \MarioBrosName, \MsPacManName, and \SurroundName, consistent with the increase of \gls{acu} with depth observed for these environments in Figure~\ref{fig:latent_space_curves_by_layer_depth}, whereas it is an intermediate layer for \BreakoutName, \PongName, and \SpaceInvadersName.
\PongName reaches the highest \gls{acu} ($0.45 \pm 0.07$) among the six environments, combining a moderate \gls{cu} with an \gls{as} that saturates from layer 3 onward, while \MarioBrosName combines the highest \gls{as} ($0.94 \pm 0.02$) with the lowest \gls{cu} ($0.37 \pm 0.08$), yielding one of the lowest \gls{acu} values.

\FloatBarrier
